\chapter{Giới thiệu (Introduction)}

\section{Bối cảnh vấn đề}
Trong những năm gần đây, sự phát triển mạnh mẽ của các nghiên cứu khoa học đã tạo ra một sự bùng nổ về khối lượng tài liệu học thuật. Các kho lưu trữ mở như arXiv hiện đang chứa hàng triệu bài báo trải dài trên hàng trăm lĩnh vực khác nhau, tạo thành một nguồn dữ liệu khổng lồ (Big Data). Mặc dù sự phong phú này là một tài nguyên quý giá, nhưng nó cũng trực tiếp dẫn đến bài toán "quá tải thông tin". Sinh viên và các nhà nghiên cứu đang phải dành một lượng thời gian khổng lồ chỉ để tra cứu, sàng lọc và đánh giá mức độ liên quan của tài liệu, thay vì tập trung vào việc nghiên cứu cốt lõi.\\
Đứng trước khối lượng dữ liệu khổng lồ này, các công cụ tìm kiếm học thuật truyền thống bắt đầu bộc lộ sự quá tải và những giới hạn về mặt kiến trúc.\\
\begin{itemize}
    \item Nếu áp dụng các công cụ tìm kiếm thuần dựa trên từ khóa (Lexical Search), hệ thống sẽ vô cùng cứng nhắc, không thể hiểu được ngữ cảnh phức tạp, từ đồng nghĩa hay ý đồ thực sự của người tra cứu.
    \item Ngược lại, nếu chỉ lạm dụng các mô hình Trí tuệ nhân tạo (AI Search) để hiểu ngữ nghĩa, hệ thống lại dễ dàng đánh mất tính chính xác tuyệt đối—một yếu tố sống còn trong nghiên cứu khoa học khi người dùng cần tìm đích danh một tác giả cụ thể, một công thức, hay một thuật ngữ viết tắt chuyên ngành hiếm gặp. Hơn nữa, việc áp dụng các bộ lọc siêu dữ liệu (như năm xuất bản, danh mục) lên các hệ thống AI thường xuyên gây ra sự sai lệch hoặc làm giảm hiệu suất hệ thống đáng kể.
\end{itemize}
Với lượng dữ liệu lớn nhưng thiếu công cụ trích xuất thông minh đã đặt ra một yêu cầu cấp thiết: Cần phải có một kiến trúc tìm kiếm thế hệ mới, có khả năng xử lý Dữ liệu lớn với tốc độ cao, đồng thời dung hòa được cả độ chính xác của tìm kiếm từ khóa và sự thông minh của AI hiểu ngữ cảnh \cite{manning2008ir}.
\section{Mục tiêu dự án}
Dự án hướng tới việc xây dựng một hệ thống tìm kiếm lai (Hybrid Search) toàn diện dành cho tài liệu học thuật, sử dụng Elasticsearch làm nền tảng cốt lõi. Bằng cách kết hợp ưu điểm của cả hai phương pháp: tìm kiếm văn bản truyền thống (Keyword Search) để đảm bảo độ chính xác tuyệt đối của các thuật ngữ chuyên ngành, và tìm kiếm vector (Vector Search) để nắm bắt ý nghĩa ngữ cảnh, hệ thống kỳ vọng sẽ khắc phục triệt để các hạn chế của các bộ máy tìm kiếm hiện tại, từ đó mang lại kết quả tra cứu chính xác và tối ưu nhất cho người dùng.
\begin{itemize}
    \item \textbf{Về quản trị và xử lý dữ liệu:} Thiết lập quy trình thu thập, làm sạch và chuẩn hóa khối lượng lớn dữ liệu học thuật thô từ các nguồn mở. Đảm bảo dữ liệu được định dạng tối ưu, sẵn sàng cho việc lập chỉ mục (indexing) và truy xuất ở quy mô hàng triệu bản ghi.
    \item \textbf{Về kiến trúc nền tảng:} Triển khai một hạ tầng tìm kiếm phân tán mạnh mẽ, có khả năng lưu trữ và xử lý đồng thời cả dữ liệu văn bản thuần túy và các biểu diễn không gian vector đa chiều. Hệ thống phải đảm bảo tính sẵn sàng cao và phản hồi truy vấn với độ trễ thấp (low latency).
    \item \textbf{Về thuật toán truy xuất cốt lõi:} Phát triển cơ chế tìm kiếm đa luồng, cho phép thực thi song song các truy vấn dựa trên từ vựng và truy vấn dựa trên ngữ nghĩa. Trọng tâm là xây dựng cấu trúc tự động đánh giá, tính điểm và hợp nhất xếp hạng từ các nguồn kết quả khác nhau để tối ưu hóa độ chính xác.
    \item \textbf{Về tính ứng dụng và đánh giá:} Hoàn thiện giao diện tương tác trực quan, hỗ trợ người dùng thực hiện các thao tác tra cứu và lọc dữ liệu phức tạp. Đồng thời, áp dụng các bộ tiêu chí đánh giá định lượng khoa học để chứng minh hiệu năng của mô hình Tìm kiếm Lai so với các phương pháp truyền thống.
\end{itemize}
\section{Phạm vi (Scope)}
Nhằm tối ưu hóa hiệu quả thực hiện và bám sát các mục tiêu kỹ thuật cốt lõi, phạm vi nghiên cứu của dự án được giới hạn trong các phân đoạn cụ thể sau:
\begin{itemize}
    \item \textbf{Về nguồn tài nguyên dữ liệu}: Dự án khai thác kho lưu trữ tài liệu học thuật mở arXiv thông qua bộ dữ liệu từ Kaggle. Đây là nguồn tài nguyên tri thức đa ngành quy mô lớn, cung cấp hệ thống thuộc tính phong phú từ nội dung văn bản cho đến các siêu dữ liệu phân loại, tạo cơ sở vững chắc cho việc thực nghiệm các mô hình truy xuất thông tin phức tạp.
    \item \textbf{Về quy mô triển khai giai đoạn giữa kỳ}: Nhóm thực hiện khoanh vùng và xử lý trên một tập dữ liệu con tiêu biểu gồm 100.000 bài báo thuộc lĩnh vực Khoa học máy tính. Việc giới hạn quy mô trong một chuyên ngành cụ thể giúp tối ưu hóa quy trình làm sạch dữ liệu, chuẩn hóa các thuật ngữ chuyên biệt và đảm bảo tính kiểm soát cao trong các kịch bản đánh giá độ chính xác của kết quả tìm kiếm.
    \item \textbf{Về hạ tầng và kỹ thuật thực hiện}: Phạm vi triển khai hạ tầng tập trung vào việc thiết lập hệ thống tìm kiếm phân tán trên hệ sinh thái Elasticsearch phiên bản 8.x. Các cấu hình lõi bao gồm định nghĩa cấu trúc chỉ mục, thiết lập các bộ phân tích từ vựng chuyên sâu và tối ưu hóa hiệu suất nạp dữ liệu quy mô lớn thông qua các phương thức xử lý hàng loạt.
    \item \textbf{Về phương pháp luận và chức năng cốt lõi}: Trọng tâm của giai đoạn hiện tại là hoàn thiện hệ thống tìm kiếm từ vựng dựa trên thuật toán BM25. Hệ thống này được xác lập như một hệ thống cơ sở để làm tiền đề quan trọng giúp nhóm tiến hành đối chứng và đo lường mức độ cải thiện hiệu năng khi tích hợp các công nghệ tìm kiếm nâng cao trong các giai đoạn tiếp theo của đồ án.
\end{itemize}